\FigureLayoutDeclare{img/2017/zhejiang/q7_fig1.png}{5.1cm}{75bp 48bp 50bp 0bp}

\chapter{2017年浙江卷}

\section{单选题}

\begin{problem}
已知集合 \(P=\{x\mid -1<x<1\}\)，\(Q=\{x\mid 0<x<2\}\)，那么 \(P\cup Q=\)（\quad）

\choices
    {\(\left(-1,2\right)\)}
    {\(\left(0,1\right)\)}
    {\(\left(-1,0\right)\)}
    {\(\left(1,2\right)\)}
\end{problem}

\begin{answer}
A
\end{answer}

\begin{solution}
由并集的定义，\(P\cup Q\) 包含两个集合中的所有元素. 区间 \(P=(-1,1)\) 与 \(Q=(0,2)\) 有公共部分，合并后覆盖从 \(-1\) 到 \(2\) 之间的所有数，且两个端点均不取，因此
\[
P\cup Q=(-1,2)
\]
故选 A.
\end{solution}

\begin{problem}
椭圆 \(\dfrac{x^{2}}{9}+\dfrac{y^{2}}{4}=1\) 的离心率是（\quad）

\choices
    {\(\dfrac{\sqrt{13}}{3}\)}
    {\(\dfrac{\sqrt5}{3}\)}
    {\(\dfrac23\)}
    {\(\dfrac59\)}
\end{problem}

\begin{answer}
B
\end{answer}

\begin{solution}
椭圆的标准方程中 \(a^2=9\)，\(b^2=4\)，所以 \(a=3\). 由 \(c^2=a^2-b^2\)，得 \(c^2=9-4=5\)，从而 \(c=\sqrt{5}\). 因此离心率为
\[
e=\frac{c}{a}=\frac{\sqrt{5}}{3}
\]
故选 B.
\end{solution}

\begin{problem}
某几何体的三视图如图所示（单位：\(\mathrm{cm}\)），则该几何体的体积（单位：\(\mathrm{cm}^3\)）是（\quad）

\begin{center}
\bitmapfigure[max width=0.30\textwidth]{img/2017/zhejiang/q3_fig1.png}
\end{center}

\choices
    {\(\dfrac{\pi}{2}+1\)}
    {\(\dfrac{\pi}{2}+3\)}
    {\(\dfrac{3\pi}{2}+1\)}
    {\(\dfrac{3\pi}{2}+3\)}
\end{problem}

\begin{answer}
A
\end{answer}

\begin{solution}
由三视图可知，该几何体可分为一个底面半径为 \(1\)、高为 \(3\) 的半圆锥，以及一个底面为等腰直角三角形、高为 \(3\) 的三棱锥. 等腰直角三角形的斜边为 \(2\)，故其面积为 \(1\).

半圆锥的体积为
\[
\frac12\cdot\frac13\cdot\pi\cdot1^2\cdot3=\frac{\pi}{2}
\]
三棱锥的体积为
\[
\frac13\cdot1\cdot3=1
\]
所以该几何体的体积为
\[
V=\frac{\pi}{2}+1
\]
故选 A.
\end{solution}

\begin{problem}
若 \(x\)，\(y\) 满足约束条件
\(\begin{cases}
x\ge0,\\
x+y-3\ge0,\\
x-2y\le0,
\end{cases}\)
则 \(z=x+2y\) 的取值范围是（\quad）

\choices
    {\([0,6]\)}
    {\([0,4]\)}
    {\([6,+\infty)\)}
    {\([4,+\infty)\)}
\end{problem}

\begin{answer}
D
\end{answer}

\begin{solution}
约束条件等价于
\(x\geq0\)，\(y\geq3-x\)，\(y\geq\frac{x}{2}\).
两条边界直线 \(y=3-x\) 与 \(y=x/2\) 的交点为 \((2,1)\). 当 \(0\leq x\leq2\) 时，下界为 \(y=3-x\)，于是
\[
z=x+2y\geq x+2(3-x)=6-x\geq4
\]
当 \(x\geq2\) 时，下界为 \(y=x/2\)，于是
\[
z=x+2y\geq x+2\cdot\frac{x}{2}=2x\geq4
\]
在 \((2,1)\) 处取到 \(z=2+2=4\)，而令 \(y\) 无限增大时 \(z\) 也无限增大. 因此
\[
z\in[4,+\infty)
\]
故选 D.
\end{solution}

\begin{problem}
若函数 \(f(x)=x^{2}+ax+b\) 在区间 \([0,1]\) 上的最大值是 \(M\)，最小值是 \(m\)，则 \(M-m\)（\quad）

\choices
    {与 \(a\) 有关，且与 \(b\) 有关}
    {与 \(a\) 有关，但与 \(b\) 无关}
    {与 \(a\) 无关，且与 \(b\) 无关}
    {与 \(a\) 无关，但与 \(b\) 有关}
\end{problem}

\begin{answer}
B
\end{answer}

\begin{solution}
令 \(g(x)=x^2+ax\)，则 \(f(x)=g(x)+b\). 在区间 \([0,1]\) 上，给 \(g\) 的所有函数值同时加上 \(b\)，只会使最大值和最小值分别都增加 \(b\)，所以 \(M-m\) 与 \(b\) 无关.

但 \(M-m\) 与 \(a\) 有关. 例如 \(a=0\) 时，\(g(x)=x^2\)，其在 \([0,1]\) 上的值域为 \([0,1]\)，故 \(M-m=1\)；而 \(a=-1\) 时，\(g(x)=x^2-x\)，其最小值为 \(-1/4\)，最大值为 \(0\)，故 \(M-m=1/4\). 因此 \(M-m\) 与 \(a\) 有关、与 \(b\) 无关.
\end{solution}

\begin{problem}
已知等差数列 \(\{a_n\}\) 的公差为 \(d\)，前 \(n\) 项和为 \(S_n\)，则 “\(d>0\)” 是 “\(S_4+S_6>2S_5\)” 的（\quad）

\choices
    {充分不必要条件}
    {必要不充分条件}
    {充分必要条件}
    {既不充分也不必要条件}
\end{problem}

\begin{answer}
C
\end{answer}

\begin{solution}
等差数列的前 \(n\) 项和为
\[
S_n=na_1+\frac{n(n-1)}{2}d
\]
因此
\(S_4=4a_1+6d\)，\(S_5=5a_1+10d\)，\(S_6=6a_1+15d\).
于是
\[
S_4+S_6-2S_5=(4a_1+6d)+(6a_1+15d)-2(5a_1+10d)=d
\]
所以 \(S_4+S_6>2S_5\) 当且仅当 \(d>0\). 两个条件互相推出，故“\(d>0\)”是“\(S_4+S_6>2S_5\)”的充分必要条件，选 C.
\end{solution}

\begin{problem}
函数 \(y=f(x)\) 的导函数 \(y=f'(x)\) 的图象如图所示，则函数 \(y=f(x)\) 的图象可能是（\quad）

\begin{center}
\bitmapfigure[max width=0.26\textwidth]{img/2017/zhejiang/q7_fig1.png}
\end{center}

\choices
    {\choicebitmap[width=0.15\paperwidth]{img/2017/zhejiang/q7_optA.png}}
    {\choicebitmap[width=0.15\paperwidth]{img/2017/zhejiang/q7_optB.png}}
    {\choicebitmap[width=0.15\paperwidth]{img/2017/zhejiang/q7_optC.png}}
    {\choicebitmap[width=0.15\paperwidth]{img/2017/zhejiang/q7_optD.png}}
\end{problem}

\begin{answer}
D
\end{answer}

\begin{solution}
由导函数图象可知，\(f'(x)\) 的符号依次为负、正、负、正，因此 \(f(x)\) 依次递减、递增、递减、递增. 故 \(f(x)\) 有两个极小值点和一个极大值点，且中间的极大值点位于 \(y\) 轴右侧.

四个选项中，只有 D 的图象具有上述单调性变化和极大值点位置，因此函数 \(y=f(x)\) 的图象可能是 D.
\end{solution}

\begin{problem}
已知随机变量 \(\xi_i\) 满足 \(P(\xi_i=1)=p_i\)，\(P(\xi_i=0)=1-p_i\)，\(i=1\)，\(2\). 若 \(0<p_1<p_2<\dfrac12\)，则（\quad）

\choices
    {\(E(\xi_1)<E(\xi_2)\)，\(D(\xi_1)<D(\xi_2)\)}
    {\(E(\xi_1)<E(\xi_2)\)，\(D(\xi_1)>D(\xi_2)\)}
    {\(E(\xi_1)>E(\xi_2)\)，\(D(\xi_1)<D(\xi_2)\)}
    {\(E(\xi_1)>E(\xi_2)\)，\(D(\xi_1)>D(\xi_2)\)}
\end{problem}

\begin{answer}
A
\end{answer}

\begin{solution}
因为 \(\xi_i\) 只取 \(0,1\)，所以
\[
E(\xi_i)=1\cdot p_i+0\cdot(1-p_i)=p_i
\]
从而 \(E(\xi_1)<E(\xi_2)\).

又
\[
D(\xi_i)=p_i(1-p_i)
\]
函数 \(h(p)=p(1-p)\) 在 \((0,1/2)\) 上的导数为 \(h'(p)=1-2p>0\)，所以由 \(0<p_1<p_2<1/2\) 得
\[
D(\xi_1)=h(p_1)<h(p_2)=D(\xi_2)
\]
故选 A.
\end{solution}

\begin{problem}
如图，已知正四面体 \(D-ABC\)（所有棱长均相等的三棱锥），\(P\)，\(Q\)，\(R\) 分别为 \(AB\)，\(BC\)，\(CA\) 上的点，\(AP=PB\)，\(\dfrac{BQ}{QC}=\dfrac{CR}{RA}=2\)，分别记二面角 \(D-PR-Q\)，\(D-PQ-R\)，\(D-QR-P\) 的平面角为 \(\alpha\)，\(\beta\)，\(\gamma\)，则（\quad）

\begin{center}
\bitmapfigure[max width=0.27\textwidth]{img/2017/zhejiang/q9_fig1.png}
\end{center}

\choices
    {\(\gamma<\alpha<\beta\)}
    {\(\alpha<\gamma<\beta\)}
    {\(\alpha<\beta<\gamma\)}
    {\(\beta<\gamma<\alpha\)}
\end{problem}

\begin{answer}
B
\end{answer}

\begin{solution}
取底面中心为 \(O\)，建立空间直角坐标系. 设正四面体棱长为 \(2\)，令
\(A=(-1,0,0)\)，\(B=(1,0,0)\)，\(C=(0,\sqrt{3},0)\)，\(P=(0,0,0)\).
正四面体顶点 \(D\) 在底面中心正上方，故可取
\[
D=\left(0,\frac{\sqrt{3}}{3},\frac{2\sqrt{6}}{3}\right)
\]
由 \(BQ:QC=2:1\)、\(CR:RA=2:1\)，得
\(Q=\left(\frac13,\frac{2\sqrt3}{3},0\right)\)，\(R=\left(-\frac23,\frac{\sqrt3}{3},0\right)\).
三点 \(P\)，\(Q\)，\(R\) 在底面 \(z=0\) 内，故该平面的法向量可取 \(\bs{m}=(0,0,1)\). 分别取
\[
\bs{n}_\alpha=9\,\overrightarrow{PR}\times\overrightarrow{PD}=(6\sqrt2,4\sqrt6,-2\sqrt3)
\]
\[
\bs{n}_\beta=9\,\overrightarrow{PQ}\times\overrightarrow{PD}=(12\sqrt2,-2\sqrt6,\sqrt3)
\]
\[
\bs{n}_\gamma=9\,\overrightarrow{QD}\times\overrightarrow{QR}=(6\sqrt2,-6\sqrt6,-2\sqrt3)
\]
它们分别是平面 \(DPR\)、\(DPQ\)、\(DQR\) 的法向量. 由法向量夹角公式，且三个平面角均取锐角，有
\[
\cos\alpha=\frac{|\bs{n}_\alpha\cdot\bs{m}|}{|\bs{n}_\alpha||\bs{m}|}=\frac{1}{\sqrt{15}}
\]
\(\cos\beta=\frac{|\bs{n}_\beta\cdot\bs{m}|}{|\bs{n}_\beta||\bs{m}|}=\frac{1}{\sqrt{105}}\)，\(\cos\gamma=\frac{|\bs{n}_\gamma\cdot\bs{m}|}{|\bs{n}_\gamma||\bs{m}|}=\frac15\).
因为 \(1/\sqrt{15}>1/5>1/\sqrt{105}\)，故
\[
\alpha<\gamma<\beta
\]
故选 B.
\end{solution}

\begin{problem}
如图，已知平面四边形 \(ABCD\)，\(AB\perp BC\)，\(AB=BC=AD=2\)，\(CD=3\)，\(AC\) 与 \(BD\) 交于点 \(O\)，记 \(I_1=\overrightarrow{OA}\cdot\overrightarrow{OB}\)，\(I_2=\overrightarrow{OB}\cdot\overrightarrow{OC}\)，\(I_3=\overrightarrow{OC}\cdot\overrightarrow{OD}\)，则（\quad）

\begin{center}
\bitmapfigure[max width=0.27\textwidth]{img/2017/zhejiang/q10_fig1.png}
\end{center}

\choices
    {\(I_1<I_2<I_3\)}
    {\(I_1<I_3<I_2\)}
    {\(I_3<I_1<I_2\)}
    {\(I_2<I_1<I_3\)}
\end{problem}

\begin{answer}
C
\end{answer}

\begin{solution}
建立平面直角坐标系，令 \(B=(0,0)\)，\(C=(2,0)\)，\(A=(0,2)\). 设 \(D=(u,v)\). 由 \(AD=2\)、\(CD=3\)，得
\(u^2+(v-2)^2=4\)，\((u-2)^2+v^2=9\).
由题图中四边形的取向取与 \(D\) 位于右上方相符的交点，解得
\(u=\frac{3+\sqrt{119}}8\)，\(v=\frac{13+\sqrt{119}}8\)
从而 \(0<u<v\) 且 \(u+v=2+\sqrt{119}/4>4\).

设 \(O=AC\cap BD\). 因为 \(AC\) 的方程为 \(x+y=2\)，而 \(O\) 在直线 \(BD\) 上，所以
\[
O=\left(\frac{2u}{u+v},\frac{2v}{u+v}\right)=:(r,2-r)
\]
由 \(u<v\) 得 \(0<r<1\). 于是
\(\overrightarrow{OA}=(-r,r)\)，\(\overrightarrow{OB}=(-r,r-2)\)
故
\[
I_1=\overrightarrow{OA}\cdot\overrightarrow{OB}=2r(r-1)<0
\]
又 \(O\) 在线段 \(AC\)、\(BD\) 内部，且
\(\frac{OA}{OC}=\frac{u}{v}<1\)，\(\frac{OB}{OD}=\frac{2}{u+v-2}<1\)
所以 \(OA<OC\)、\(OB<OD\). 由于 \(\overrightarrow{OC}\) 与 \(\overrightarrow{OA}\) 方向相反，\(\overrightarrow{OD}\) 与 \(\overrightarrow{OB}\) 方向相反，得到
\[
I_2=\overrightarrow{OB}\cdot\overrightarrow{OC}>0
\]
且
\[
I_3=\overrightarrow{OC}\cdot\overrightarrow{OD}
=\frac{OC\cdot OD}{OA\cdot OB}I_1<I_1
\]
其中比例因子大于 \(1\)，而 \(I_1<0\). 因此
\[
I_3<I_1<I_2
\]
故选 C.
\end{solution}

\section{填空题}

\begin{problem}
我国古代数学家刘徽创立的“割圆术”可以估算圆周率 \(\pi\)，理论上能把 \(\pi\) 的值计算到任意精度，祖冲之继承并发展了“割圆术”，将 \(\pi\) 的值精确到小数点后七位，其结果领先世界一千多年，“割圆术”的第一步是计算单位圆内接正六边形的面积 \(S_6\)，\(S_6=\fillinblank{}\).
\end{problem}

\begin{answer}
\(\dfrac{3\sqrt{3}}{2}\)
\end{answer}

\begin{solution}
单位圆的半径为 \(1\). 正六边形的每条边等于外接圆半径，故边长为 \(1\). 连接圆心与正六边形的六个顶点，将正六边形分成六个边长为 \(1\) 的正三角形，每个正三角形的面积为
\[
\frac12\cdot1\cdot1\cdot\sin60^\circ=\frac{\sqrt3}{4}
\]
所以
\[
S_6=6\cdot\frac{\sqrt3}{4}=\frac{3\sqrt3}{2}
\]
\end{solution}

\begin{problem}
已知 \(a\)，\(b\in\R\)，\((a+b\i)^2=3+4\i\)（\(\i\) 是虚数单位），则 \(a^2+b^2=\fillinblank{}\)，\(ab=\fillinblank{}\).
\end{problem}

\begin{answer}
\(5\)，\(2\)
\end{answer}

\begin{solution}
展开复数等式左边，得
\[
(a+b\i)^2=a^2-b^2+2ab\i
\]
与 \(3+4\i\) 比较实部和虚部，得到
\(a^2-b^2=3\)，\(2ab=4\).
因此 \(ab=2\). 又
\[
(a^2+b^2)^2=(a^2-b^2)^2+(2ab)^2=3^2+4^2=25
\]
由于 \(a^2+b^2\geq0\)，所以 \(a^2+b^2=5\).
\end{solution}

\begin{problem}
已知多项式 \((x+1)^3(x+2)^2=x^5+a_1x^4+a_2x^3+a_3x^2+a_4x+a_5\)，则 \(a_4=\fillinblank{}\)，\(a_5=\fillinblank{}\).
\end{problem}

\begin{answer}
\(16\)，\(4\)
\end{answer}

\begin{solution}
先展开
\((x+1)^3=x^3+3x^2+3x+1\)，\((x+2)^2=x^2+4x+4\).
其中 \(x\) 的系数来自 \(3x\cdot4\) 与 \(1\cdot4x\)，所以
\[
a_4=3\cdot4+1\cdot4=16
\]
常数项为 \(1\cdot4=4\)，所以 \(a_5=4\).
\end{solution}

\begin{problem}
已知 \(\triangle ABC\)，\(AB=AC=4\)，\(BC=2\)，点 \(D\) 为 \(AB\) 延长线上一点，\(BD=2\)，连接 \(CD\)，则 \(\triangle BDC\) 的面积是\fillinblank{}，\(\cos\angle BDC=\fillinblank{}\).
\end{problem}

\begin{answer}
\(\dfrac{\sqrt{15}}{2}\)，\(\dfrac{\sqrt{10}}{4}\)
\end{answer}

\begin{solution}
由 \(AB=AC=4\)、\(BC=2\)，在等腰三角形 \(ABC\) 中
\[
\cos\angle ABC=\frac{AB^2+BC^2-AC^2}{2\cdot AB\cdot BC}=\frac14
\]
故 \(\sin\angle ABC=\sqrt{1-1/16}=\sqrt{15}/4\). 由于 \(BD=2<AB=4\)，点 \(D\) 只能在 \(B\) 的外侧，故射线 \(BD\) 与射线 \(BA\) 反向，所以
\[
\sin\angle DBC=\sin\angle ABC=\frac{\sqrt{15}}4
\]
于是
\[
S_{\triangle BDC}=\frac12\cdot BD\cdot BC\cdot\sin\angle DBC
=\frac12\cdot2\cdot2\cdot\frac{\sqrt{15}}4=\frac{\sqrt{15}}2
\]
又 \(\cos\angle DBC=-1/4\)，由余弦定理
\[
CD^2=BD^2+BC^2-2\cdot BD\cdot BC\cos\angle DBC
=4+4-8\left(-\frac14\right)=10
\]
在三角形 \(BDC\) 中再次用余弦定理，得
\[
\cos\angle BDC=\frac{BD^2+CD^2-BC^2}{2\cdot BD\cdot CD}
=\frac{4+10-4}{4\sqrt{10}}=\frac{\sqrt{10}}4
\]
\end{solution}

\begin{problem}
已知向量 \(\bs{a}\)，\(\bs{b}\) 满足 \(|\bs{a}|=1\)，\(|\bs{b}|=2\)，则 \(|\bs{a}+\bs{b}|+|\bs{a}-\bs{b}|\) 的最小值是\fillinblank{}，最大值是\fillinblank{}.
\end{problem}

\begin{answer}
\(4\)，\(2\sqrt{5}\)
\end{answer}

\begin{solution}
设 \(\theta\) 为 \(\bs a\) 与 \(\bs b\) 的夹角，则
\[
|\bs a+\bs b|^2=1+4+4\cos\theta=5+4\cos\theta
\]
\[
|\bs a-\bs b|^2=1+4-4\cos\theta=5-4\cos\theta
\]
记两者之和为 \(T\). 则
\[
T^2=10+2\sqrt{(5+4\cos\theta)(5-4\cos\theta)}
=10+2\sqrt{25-16\cos^2\theta}
\]
因为 \(0\leq\cos^2\theta\leq1\)，所以
\[
16\leq T^2\leq20
\]
当 \(|\cos\theta|=1\) 时取下界，\(T=4\)；当 \(\cos\theta=0\) 时取上界，\(T=2\sqrt5\). 因而最小值为 \(4\)，最大值为 \(2\sqrt5\).
\end{solution}

\begin{problem}
从 \(6\) 男 \(2\) 女共 \(8\) 名学生中选出队长 \(1\) 人，副队长 \(1\) 人，普通队员 \(2\) 人组成 \(4\) 人服务队，要求服务队中至少有 \(1\) 名女生，共有\fillinblank{} 种不同的选法.
\end{problem}

\begin{answer}
\(660\)
\end{answer}

\begin{solution}
按服务队中女生人数分类计数. 恰有 \(1\) 名女生时，先选出该女生和 \(3\) 名男生，再安排队长、副队长，普通队员的两个位置不区分，选法数为
\[
\binom21\binom63\cdot4\cdot3=480
\]
恰有 \(2\) 名女生时，选出 \(2\) 名女生和 \(2\) 名男生，再安排队长和副队长，选法数为
\[
\binom22\binom62\cdot4\cdot3=180
\]
所以符合条件的选法共有
\[
480+180=660
\]
\end{solution}

\begin{problem}
已知 \(a\in\R\)，函数 \(f(x)=\left|x+\dfrac{4}{x}-a\right|+a\) 在区间 \([1,4]\) 上的最大值是 5，则 \(a\) 的取值范围是\fillinblank{}.
\end{problem}

\begin{answer}
\(\left(-\infty,\dfrac{9}{2}\right]\)
\end{answer}

\begin{solution}
令 \(g(x)=x+4/x\). 则
\[
g'(x)=1-\frac4{x^2}
\]
所以 \(g\) 在 \([1,2]\) 上递减，在 \([2,4]\) 上递增，并且
\(g(2)=4\)，\(g(1)=g(4)=5\).
故 \(g([1,4])=[4,5]\). 因此
\[
\max_{x\in[1,4]}f(x)=a+\max_{t\in[4,5]}|t-a|
\]
当 \(a\leq9/2\) 时，区间 \([4,5]\) 中距 \(a\) 最远的端点为 \(5\)，最大值为
\[
a+|5-a|=a+5-a=5
\]
当 \(a>9/2\) 时，若 \(a\leq5\)，最大值为 \(a+(a-4)=2a-4>5\)；若 \(a>5\)，最大值仍为 \(a+(a-4)=2a-4>5\). 因而满足条件的取值范围为
\[
a\in\left(-\infty,\frac92\right]
\]
\end{solution}

\section{解答题}

\begin{problem}
已知函数 \(f(x)=\sin^2x-\cos^2x-2\sqrt3\sin x\cos x\)（\(x\in\R\)）.

\begin{enumerate}
\item 求 \(f\left(\dfrac{2\pi}{3}\right)\) 的值.
\item 求 \(f(x)\) 的最小正周期及单调递增区间.
\end{enumerate}
\end{problem}

\begin{solution}
\begin{enumerate}
\item 由二倍角公式，\(f(x)=-\cos 2x-\sqrt{3}\sin 2x=2\sin\left(-2x-\dfrac{\pi}{6}\right)\)，所以
\(f\left(\dfrac{2\pi}{3}\right)=2\sin\left(-\dfrac{3\pi}{2}\right)=2\).
\item 由 \(f(x)=2\sin\left(-2x-\dfrac{\pi}{6}\right)\) 可知，\(f(x)\) 的最小正周期为 \(\pi\). 又 \(f'(x)=-4\cos\left(2x+\dfrac{\pi}{6}\right)\)，因此当 \(\dfrac{\pi}{2}+2k\pi\leq 2x+\dfrac{\pi}{6}\leq\dfrac{3\pi}{2}+2k\pi\) 时，\(f'(x)\geq0\)，其中 \(k\in\Z\). 解得单调递增区间为 \(\left[\dfrac{\pi}{6}+k\pi,\dfrac{2\pi}{3}+k\pi\right]\)，\(k\in\Z\).
\end{enumerate}
\end{solution}

\begin{problem}
如图，已知四棱锥 \(P-ABCD\)，\(\triangle PAD\) 是以 \(AD\) 为斜边的等腰直角三角形，\(BC\parallel AD\)，\(CD\perp AD\)，\(PC=AD=2DC=2CB\)，\(E\) 为 \(PD\) 的中点.

\begin{enumerate}
\item 证明：\(CE\parallel\) 平面 \(PAB\)；
\item 求直线 \(CE\) 与平面 \(PBC\) 所成角的正弦值.
\end{enumerate}

\begin{center}
\bitmapfigure[max width=0.30\textwidth]{img/2017/zhejiang/q19_fig1.png}
\end{center}
\end{problem}

\begin{solution}
\begin{enumerate}
\item 取 \(AD\) 的中点 \(F\)，连接 \(EF\)，\(CF\). 因为 \(E\)，\(F\) 分别为 \(PD\)，\(AD\) 的中点，所以 \(EF\parallel PA\). 又因为 \(BC\parallel AD\)，且 \(BC=\dfrac{1}{2}AD\)，由 \(F\) 为 \(AD\) 的中点可得 \(CF\parallel AB\). 因此平面 \(EFC\parallel\) 平面 \(PAB\)，而 \(CE\subset\) 平面 \(EFC\)，所以 \(CE\parallel\) 平面 \(PAB\).
\item 设 \(PC=AD=2DC=2CB=2\)，以 \(D\) 为原点，\(DA\)，\(DC\) 所在直线分别为 \(x\)，\(y\) 轴，过 \(D\) 作底面 \(ABCD\) 的垂线为 \(z\) 轴，建立空间直角坐标系，则 \(A(2,0,0)\)，\(C(0,1,0)\)，\(B(1,1,0)\). 设 \(P(1,u,v)\)，由 \(PA=PD\) 得 \(P\) 的横坐标为 \(1\)，由 \(PD=\sqrt{2}\) 得 \(u^2+v^2=1\)，再由 \(PC=2\) 得 \(1+(u-1)^2+v^2=4\)，所以 \(u=-\dfrac{1}{2}\)，\(v=\dfrac{\sqrt{3}}{2}\). 于是 \(P\left(1,-\dfrac{1}{2},\dfrac{\sqrt{3}}{2}\right)\)，\(E\left(\dfrac{1}{2},-\dfrac{1}{4},\dfrac{\sqrt{3}}{4}\right)\)，从而 \(\overrightarrow{CE}=\left(\dfrac{1}{2},-\dfrac{5}{4},\dfrac{\sqrt{3}}{4}\right)\). 取平面 \(PBC\) 的一个法向量 \(\bs{n}=\overrightarrow{BC}\times\overrightarrow{BP}=(0,\sqrt{3},3)\)，则 \(\overrightarrow{CE}\cdot\bs{n}=-\dfrac{\sqrt{3}}{2}\)，\(|\overrightarrow{CE}|=\sqrt{2}\)，\(|\bs{n}|=2\sqrt{3}\). 设直线 \(CE\) 与平面 \(PBC\) 所成的角为 \(\theta\)，则 \(\sin\theta=\dfrac{|\overrightarrow{CE}\cdot\bs{n}|}{|\overrightarrow{CE}|\,|\bs{n}|}=\dfrac{\sqrt{3}/2}{\sqrt{2}\cdot2\sqrt{3}}=\dfrac{\sqrt{2}}{8}\).
\end{enumerate}
\end{solution}

\begin{problem}
已知函数 \(f(x)=\bigl(x-\sqrt{2x-1}\bigr)\e^{-x}\)（\(x\ge\dfrac12\)）.

\begin{enumerate}
\item 求 \(f(x)\) 的导函数；
\item 求 \(f(x)\) 在区间 \(\left[\dfrac12,+\infty\right)\) 上的取值范围.
\end{enumerate}
\end{problem}

\begin{solution}
\begin{enumerate}
\item 当 \(x>\dfrac{1}{2}\) 时，利用乘积法则得 \(f'(x)=\left(1-\dfrac{1}{\sqrt{2x-1}}\right)\e^{-x}-\left(x-\sqrt{2x-1}\right)\e^{-x}=\left(1-x\right)\left(1-\dfrac{2}{\sqrt{2x-1}}\right)\e^{-x}\). 在 \(x=\dfrac{1}{2}\) 处，右导数趋于 \(-\infty\)，故不存在有限导数.
\item 因为 \(\e^{-x}>0\)，所以在 \(\left(\dfrac{1}{2},1\right)\) 上 \(f'(x)<0\)，在 \(\left(1,\dfrac{5}{2}\right)\) 上 \(f'(x)>0\)，在 \(\left(\dfrac{5}{2},+\infty\right)\) 上 \(f'(x)<0\). 又 \(f\left(\dfrac{1}{2}\right)=\dfrac{1}{2}\e^{-\frac{1}{2}}\)，\(f(1)=0\)，\(f\left(\dfrac{5}{2}\right)=\dfrac{1}{2}\e^{-\frac{5}{2}}\)，且 \(\displaystyle\lim_{x\to+\infty}f(x)=0\). 此外，\(x-\sqrt{2x-1}=\dfrac{(x-1)^2}{x+\sqrt{2x-1}}\geq0\)，故 \(f(x)\geq0\). 因此 \(f(x)\) 在 \(\left[\dfrac{1}{2},+\infty\right)\) 上的取值范围为 \(\left[0,\dfrac{1}{2}\e^{-\frac{1}{2}}\right]\).
\end{enumerate}
\end{solution}

\begin{problem}
如图，已知抛物线 \(x^2=y\)，点 \(A\left(-\dfrac12,\dfrac14\right)\)，\(B\left(\dfrac32,\dfrac94\right)\)，抛物线上的点 \(P(x,y)\)（\(-\dfrac12<x<\dfrac32\)），过点 \(B\) 作直线 \(AP\) 的垂线，垂足为 \(Q\).

\begin{enumerate}
\item 求直线 \(AP\) 斜率的取值范围；
\item 求 \(|PA|\cdot|PQ|\) 的最大值.
\end{enumerate}

\begin{center}
\bitmapfigure[max width=0.30\textwidth]{img/2017/zhejiang/q21_fig1.png}
\end{center}
\end{problem}

\begin{solution}
\begin{enumerate}
\item 设 \(P(t,t^2)\)，则 \(-\dfrac{1}{2}<t<\dfrac{3}{2}\)，所以直线 \(AP\) 的斜率为 \(\dfrac{t^2-\frac{1}{4}}{t+\frac{1}{2}}=t-\dfrac{1}{2}\)，其取值范围为 \((-1,1)\).
\item 记直线 \(AP\) 的斜率为 \(k\)，则 \(t=k+\dfrac{1}{2}\)，且 \(-1<k<1\). 由 \(\overrightarrow{AP}=(1+k)(1,k)\) 可得 \(|PA|=(1+k)\sqrt{1+k^2}\). 又 \(\overrightarrow{BP}=\left(-1+k, -2+k+k^2\right)\)，取直线 \(AP\) 的方向向量 \((1,k)\)，则 \(|PQ|=\dfrac{|\overrightarrow{BP}\cdot(1,k)|}{\sqrt{1+k^2}}=\dfrac{(1-k)(1+k)^2}{\sqrt{1+k^2}}\). 因此 \(|PA|\cdot|PQ|=(1+k)^3(1-k)\). 令 \(h(k)=(1+k)^3(1-k)\)，则 \(h'(k)=2(1+k)^2(1-2k)\). 所以 \(h(k)\) 在 \((-1,\dfrac{1}{2})\) 上递增，在 \((\dfrac{1}{2},1)\) 上递减，故最大值为 \(h\left(\dfrac{1}{2}\right)=\left(\dfrac{3}{2}\right)^3\cdot\dfrac{1}{2}=\dfrac{27}{16}\).
\end{enumerate}
\end{solution}

\begin{problem}
已知数列 \(\{x_n\}\) 满足：\(x_1=1\)，\(x_n=x_{n+1}+\ln(1+x_{n+1})\)（\(n\in \N^*\)），证明：当 \(n\in \N^*\) 时，

\begin{enumerate}
\item \(0<x_{n+1}<x_n\)；
\item \(2x_{n+1}-x_n\le \dfrac{x_nx_{n+1}}{2}\)；
\item \(\dfrac{1}{2^{n-1}}\le x_n\le \dfrac{1}{2^{n-2}}\).
\end{enumerate}
\end{problem}

\begin{solution}
\begin{enumerate}
\item 令 \(\varphi(t)=t+\ln(1+t)\)，其中 \(t>-1\). 因为 \(\varphi'(t)=1+\dfrac{1}{1+t}>0\)，且 \(\varphi(0)=0\)，所以 \(\varphi\) 在定义域上单调递增. 由 \(x_1=1>0\) 归纳可知，若 \(x_n>0\)，则 \(x_n=\varphi(x_{n+1})>\varphi(0)=0\)，从而 \(x_{n+1}>0\). 于是 \(\ln(1+x_{n+1})>0\)，由递推关系得 \(x_n=x_{n+1}+\ln(1+x_{n+1})>x_{n+1}\)，所以 \(0<x_{n+1}<x_n\).
\item 令 \(t=x_{n+1}>0\)，则 \(x_n=t+\ln(1+t)\). 考虑函数 \(g(t)=t^2-2t+(t+2)\ln(1+t)\)，有 \(g(0)=0\)，\(g'(t)=2t-1+\ln(1+t)+\dfrac{1}{1+t}\)，\(g''(t)=2+\dfrac{t}{(1+t)^2}>0\). 因此 \(g'(t)>g'(0)=0\)，进而 \(g(t)>g(0)=0\)（\(t>0\)）. 代入 \(x_n=t+\ln(1+t)\)，得 \(x_nx_{n+1}-4x_{n+1}+2x_n=g(t)\geq0\)，即 \(2x_{n+1}-x_n\leq\dfrac{x_nx_{n+1}}{2}\).
\item 对任意 \(t>0\)，有 \(0<\ln(1+t)\leq t\)，所以 \(x_n\leq2x_{n+1}\)，从而 \(x_{n+1}\geq\dfrac{x_n}{2}\). 反复使用并结合 \(x_1=1\)，得到 \(x_n\geq\dfrac{1}{2^{n-1}}\). 另一方面，由第（2）问及 \(x_nx_{n+1}>0\)，得 \(\dfrac{1}{x_{n+1}}-\dfrac{1}{2}\geq2\left(\dfrac{1}{x_n}-\dfrac{1}{2}\right)\). 反复使用此不等式，得到 \(\dfrac{1}{x_n}-\dfrac{1}{2}\geq2^{n-1}\left(\dfrac{1}{x_1}-\dfrac{1}{2}\right)=2^{n-2}\)，所以 \(\dfrac{1}{x_n}\geq2^{n-2}+\dfrac{1}{2}\geq2^{n-2}\)，即 \(x_n\leq\dfrac{1}{2^{n-2}}\). 综上，\(\dfrac{1}{2^{n-1}}\leq x_n\leq\dfrac{1}{2^{n-2}}\).
\end{enumerate}
\end{solution}
